% Methods proposed by the product, B1 to B5 (the details of B1 and B2 live in their own manuscript).

\subsection{Device-aware weighted GS}\label{sec:b1}

The first proposed method puts the pixel crosstalk and the quantizer of the SLM inside the loop of the weighted GS for trap arrays, with the phase update replaced by a few conjugate-gradient steps on a device-aware loss. The hypothesis is that this makes the traps equal on the device and not only in the ideal model: on the crosstalk device the relative standard deviation of the trap intensities is at least two times below that of phase-fixed weighted GS, at an efficiency at most ten percent lower. The gap is stated against the literature. Schroff et al.~\cite{schroff} put the crosstalk model into conjugate gradient on shaped potentials, with a long time per hologram and no weighted-GS comparison; Persson, Engstr\"om and Goks\"or~\cite{persson} modify a commonly used hologram algorithm to compensate the pixel crosstalk of the SLM and improve the spot uniformity by more than one hundred percent on a device with large crosstalk; Kim et al.~\cite{kim2019} rely on camera feedback and use no device model in the solver; Ebadi et al.~\cite{ebadi2021} pre-compensate the diffraction envelope of the array and correct the aberrations of the system from measured atom light shifts, and their text describes no model of the pixel crosstalk or the quantizer. The proposal combines two published pieces, the weights of Di Leonardo et al.~\cite{dileonardo} and the device model of Schroff et al., and answers the problem differently from phase-fixed weighted GS, because the solver is aware of the device, so it is a method and not a validation protocol. The full derivation, the experiments on the factor per geometry and the ablation are in the companion manuscript~\cite{companion}; this section gives the method as implemented.

\textbf{Method.} The weights of the weighted GS are kept, $w_m\leftarrow w_m\langle r\rangle/r_m$, with $r_m=|V_m|/a_m$ the trap amplitude relative to its design amplitude, measured on the focal field of the device. The projection of the phase update is replaced by $K$ conjugate-gradient steps on the device-aware loss
\begin{equation}
L(\phi)=\frac{\sum_m\big(|V_m(\phi)|^2-T_m\big)^2}{\sum_mT_m^2},\qquad
T_m=\eta\,\frac{(w_ma_m)^2}{\sum_k(w_ka_k)^2},\qquad \eta=\sum_m|V_m|^2,
\label{eq:b1}
\end{equation}
where $V=\mathrm{Op}(\phi)$ runs through the straight-through quantizer, the crosstalk operator (the $3\times3$ upsampling and the circular convolution with the kernel of Eq.~\eqref{eq:m13}), the Fourier transform of the upsampled grid and the pixel-aperture envelope. The gradient is the closed form of the adjoint of Eq.~\eqref{eq:m13-adjoint}, so there is no automatic differentiation. The target level is the current efficiency $\eta$, so the steps flatten the traps without trading efficiency away. The weights are meant to supply the integral feedback that the few steps per iteration lack, but the ablation of Table~\ref{tab:b1-weights} does not show that they are needed: with every weight kept at one the spread is lower on the acceptance task and on the square lattice of the C08 device, and higher on the ring, so that neither form dominates and the method keeps its weights as defined. Each outer iteration therefore reads
\begin{equation}
w_m\leftarrow w_m\frac{\langle r\rangle}{r_m},\qquad
\phi\leftarrow\mathrm{CG}_K\big[L(\cdot;w);\phi\big],
\label{eq:b1-outer}
\end{equation}
restarting the conjugate gradient at every outer iteration because the target changes with the weights.

\textbf{Safeguard.} The weight update is an integrator, and three conjugate-gradient steps per iteration may not follow it on a device whose trap response is quantized and blurred. On the square lattice of the benchmark on the device of case C08 one seed in two ran away inside the model of the device-aware weighted GS itself: the loss rose by an order of magnitude between outer iterations two and four, the weights spread to a ratio of about a thousand and the efficiency collapsed. The method now returns its best iterate, the warm start included: the iterate of lowest relative standard deviation on its device model among those that keep a fraction \bOneFloor{} of the warm start's efficiency, the ten percent of the criterion. A run that converges returns its last iterate or one of its tail, and a run that winds up returns an admissible iterate no worse than its start; both are tested.

\textbf{Pseudocode.} Algorithm~\ref{alg:b1} lists the method with its safeguard.

\begin{algorithm}[t]
\caption{Device-aware weighted GS: the weight update with conjugate-gradient steps on the device-aware loss, and the safeguard that returns the best admissible iterate.}
\label{alg:b1}
\begin{algorithmic}[1]
\Require traps on bins, device, warm iterations $W$, steps per outer iteration $K_c$, outer budget $K$, option \texttt{weights}
\State $\phi\gets$ $W$ iterations of weighted GS from the random superposition (the ideal-model start); $w_m\gets1$; $a_m\gets\sqrt{w^{\mathrm{target}}_m/\langle w^{\mathrm{target}}\rangle}$
\State $V\gets$ trap fields of $\phi$ on the device; $\mathrm{floor}\gets0.9\times$ efficiency of $V$; $\mathrm{best}\gets(\phi,\sigma(V))$
\For{$i=1,\dots,K$}
  \State $V\gets$ trap fields of $\phi$ on the device; \textbf{if} \texttt{weights} \textbf{then} $r_m\gets|V_m|/a_m$; $w_m\gets w_m\langle r\rangle/r_m$ \textbf{end if}
  \State $T_m\gets\eta\,(w_ma_m)^2/\sum_k(w_ka_k)^2$ with $\eta=\sum_m|V_m|^2$
  \State $\phi\gets$ $K_c$ conjugate-gradient steps on $L(\cdot;w)$ of Eq.~\eqref{eq:b1} from $\phi$ (restarted at every outer iteration)
  \State record the history entry; \textbf{if} efficiency $\ge\mathrm{floor}$ \textbf{and} $\sigma<\sigma_{\mathrm{best}}$ \textbf{then} $\mathrm{best}\gets(\phi,\sigma)$ \textbf{end if}
\EndFor
\State \Return the phase of $\mathrm{best}$ (the warm start included)
\end{algorithmic}
\end{algorithm}

\textbf{Parameters.} The number of conjugate-gradient steps per outer iteration is \defv{devicewgs}{steps}, DinHot's choice; the warm start is \defv{devicewgs}{warm} iterations of ideal-model weighted GS (without the phase freeze; the count equals the freeze index of Kim); the outer budget of the acceptance test is four iterations. The device must have its traps on bins with no depth offset or vortex, and a shaped target is refused.

\textbf{Complexity.} Each conjugate-gradient step costs one forward and one adjoint evaluation of the upsampled model, about nine times the pixels of the plain transform; the time per iteration relative to the baseline was not measured. The method is baked.

\textbf{Acceptance criterion and gate.} On the crosstalk device the relative standard deviation is at least two times below that of phase-fixed weighted GS at an efficiency at most ten percent lower (Table~\ref{tab:ladder}). The gate takes the $10\times10$ task on the device with the published kernel, scores the weighted-GS baseline at two iteration counts, and compares the device-aware weighted GS with four outer iterations on one seed against the strongest of the two. The measurement of this report, taken from the measurements of the companion manuscript~\cite{companion}, is: \acc{devicewgs}. The report also reads the committed bake and benchmark. On the small device-limited case C08 the baked factor is \bakeCeightRatio{} on the bake's own seed, at an efficiency ratio of \bakeCeightEffRatio{}, and \bakeCeightRatioThree{} on the ratio of the means of its \bakeCeightSeedsN{} seeds. In the benchmark on the device of C08, the device-aware weighted GS meets both thresholds on \benchOnePass{} of \benchOneN{} geometries (\benchOnePassList), lowers the spread without the factor of two on the large set, and keeps its start on the lattices, where every iterate past the start gives up more than a tenth of the light; on the ideal device the spreads of the device-aware weighted GS and phase-fixed weighted GS are both below the resolution of the file on most geometries, which is a tie and not a gain. The hypothesis is therefore confirmed on the criterion of the ladder and narrowed by the benchmark.

\textbf{Limits.} The pixel of the kernel width is DinHot's reading, the SLM pixel, because the source does not settle it, and the other reading changes the width by a factor of three, on which the size of the gain depends (Section~\ref{sec:m13}). Flicker is random and is dropped from the optimization model, so it is scored but not modelled. The crosstalk parameters are device specific, because the calibration of the source takes hours; the size of the effect on a typical device of smaller pitch is unknown. The module is \texttt{b1\_device\_wgs}.

\subsection{Flicker-bounded phase-continuous sequence}\label{sec:b2}

The second proposed method concerns moving traps. The linear phase interpolation of Section~\ref{sec:m16} interpolates the position, amplitude and optical phase of every tweezer between two endpoint holograms and inverts the transform at each step, so the intermediate holograms are never individually optimized. The hypothesis is that optimizing every intermediate frame, with the trap phases locked to a smooth trajectory, keeps the minimum transient intensity of a moving trap at least 0.9 of its mean and above that of the interpolation, with a per-frame spread no worse. The gap is that the source of the interpolation observes with a high-speed camera that the weighted-GS spot nearly vanished several times, and with atoms in a separate test that phase slips near $\pi$ per step lost atoms and that slips of several tenths of $\pi$ barely mattered~\cite{lpi}; no source optimizes the frames of a trap trajectory. The experiments and the scope measurements are in the companion manuscript~\cite{companion}.

\textbf{Method.} For a trajectory $x_m(s)$, $s=0,\dots,S$, and trap phases $\theta_m(s)$ smooth in $s$, every frame is the phase-fixed weighted GS of phase-fixed weighted GS with the trap phases locked to the trajectory,
\begin{equation}
c_m=w_ma_me^{i\theta_m(s)},\qquad w_m\leftarrow w_m\frac{\langle r\rangle}{r_m},\qquad
\phi_j=\arg\sum_mc_m\,e^{+i\Delta_j^m(x_m(s))},
\label{eq:b2}
\end{equation}
started from the previous frame's hologram and weights with a few warm iterations per frame (at the first frame after the cold start the weights are one, because the cold stage does not return its own). The positions follow the minimal-jerk profile and the phases the same blend,
\begin{equation}
x_m(s)=x_m(0)+\big(x_m(S)-x_m(0)\big)h(\tau),\quad h(\tau)=10\tau^3-15\tau^4+6\tau^5,\quad \tau=\frac sS,
\label{eq:b2-traj}
\end{equation}
\begin{equation}
\theta_m(s)=\theta_m(0)+\operatorname{wrap}\big(\theta_m^{\mathrm{end}}-\theta_m(0)\big)\,h(\tau),
\label{eq:b2-phase}
\end{equation}
with the end phases held at the initial ones by default. Frame 0 is a cold phase-fixed weighted GS. The optimization replaces the penalty form of the loss, a weighted sum of squared deviations of the field phase from $\theta$, by the hard lock, its infinite-weight limit.

\textbf{Pseudocode.} Algorithm~\ref{alg:b2} lists the sequence.

\begin{algorithm}[t]
\caption{Flicker-bounded phase-continuous sequence. Each frame is a warm-started phase-fixed weighted GS whose trap phases are locked to the smooth trajectory of Eq.~\eqref{eq:b2-phase}.}
\label{alg:b2}
\begin{algorithmic}[1]
\Require initial traps, destination traps with the same identities, frame count $S+1$, frame iterations $K_f$, phase-step bound $\Delta\theta_{\max}$
\State compute $x_m(s)$ and $\theta_m(s)$ for all $m$ and $s$ by Eqs.~\eqref{eq:b2-traj} and \eqref{eq:b2-phase}
\If{$\max_{m,s}|\theta_m(s+1)-\theta_m(s)|>\Delta\theta_{\max}$} \State \textbf{raise} an error \EndIf
\State $\phi_0\gets$ cold phase-fixed weighted GS on $x(0)$; $w_m\gets1$
\For{$s=1,\dots,S$}
  \State $\phi_s,w\gets K_f$ warm iterations of Eq.~\eqref{eq:b2} at $x(s)$ and $\theta(s)$, started from $\phi_{s-1},w$
\EndFor
\State \Return frames $\phi_0,\dots,\phi_S$ and the transient ratios of Eq.~\eqref{eq:b2-ratio}
\end{algorithmic}
\end{algorithm}

\textbf{Transient intensity (DinHot's model).} Between frames $s$ and $s+1$ the displayed SLM field is modelled as the blend $f_a=A[(1-a)e^{i\phi_s}+a\,e^{i\phi_{s+1}}]$ with $a\in[0,1]$, the liquid crystal relaxing between the two holograms, and trap $m$ is read at its trajectory point $x_m(s+a)$, so that $I_m(s,a)=|V_m(f_a;x_m(s+a))|^2$. The quantity of the claim is
\begin{equation}
\rho_m=\min_{s,a}\frac{I_m(s,a)}{\operatorname{mean}_sI_m(s)},
\label{eq:b2-ratio}
\end{equation}
evaluated over a number of sub-steps per frame interval, and both the flicker-bounded sequence and the interpolation baseline of the test are scored by this same function. The source gives the measurement of the interpolation, not its transient equation, so the comparison is internally fair and not a statement about a real SLM. The refresh rate and the temporal response of the device are not modelled.

\textbf{Parameters.} The default frame count is \defv{flicker}{frames}, the length of the sort of the design study; the iterations per frame default to \defv{flicker}{frameiterations} (a mutation with none fails the acceptance test); the phase step per frame is bounded by \defv{flicker}{phasevelocitymax} rad, DinHot's bound based on the observation that slips of several tenths of $\pi$ barely mattered.

\textbf{Acceptance criterion and gate.} The minimum transient intensity of a moving trap is at least 0.9 of its mean against the interpolation (Table~\ref{tab:ladder}); the gate moves two columns of a $4\times4$ lattice by four bins in 13 frames, over three seeds, with the interpolation baseline written in the test from the text of the source, and requires, per seed, a ratio at least 0.9 and above that of the interpolation and a largest per-frame spread no worse. The measurement of this report, taken from the companion manuscript~\cite{companion}, is: \acc{flicker}. The measured scope is narrower than the criterion, as the companion manuscript reports: with the fewest frames that keep the peak step per frame below a bound set where case C10's lowest moving trap first reached 0.9, the claim holds on most seeds and not on all; on the long sort of C10 at an equal number of frames the interpolation leads; and on the seeds below the bound the limit is the phase path that the end phases prescribe, not the frame count.

\textbf{Complexity.} The sequence is $S$ warm-started runs, $O(S\,K_f\,N_{\mathrm{pix}}M)$ in the point-trap form, and the method is baked; the source of the interpolation reports a very short time per hologram for its own method on a modern graphics processor, which this product does not rely on.

\textbf{Limits.} The transient model is DinHot's: it compares sequences under one model of the display and gives no absolute intensity of a measured device; the baseline is DinHot's implementation of the published description and not the authors' code. The module is \texttt{b2\_flicker\_bounded}.

\subsection{Learned warm start for phase-fixed weighted GS}\label{sec:b3}

The third proposed method is a negative result, kept as a result. The hypothesis is that a small permutation-equivariant network over the trap list, predicting the initial trap phases $\theta_m$ of phase-fixed weighted GS, lets phase-fixed weighted GS reach a relative standard deviation of one percent in at least two times fewer iterations (median) than from a random start on held-out geometries, with the final efficiency and uniformity equal within one percent and a model of at most one megabyte. The gap is that a U-Net generates a whole phase for about twenty traps in a fraction of a second~\cite{ren2024}, DeepCGH shows learned speedups for volumes~\cite{eybposh}, and a hologram network of below 620 kilobytes is possible~\cite{shi2021} (abstract only read); no source that was read uses a network to initialize weighted GS, and none compares with it. The risk named before the work was that the random start is already reliable for the weighted GS, so that the gain may be small.

\textbf{Method.} The start is the superposition hologram with the learned trap phases,
\begin{equation}
\phi_j^{(0)}=\arg\sum_m\sqrt{w_m}\,e^{\,i(\Delta_j^m+\theta_m)},
\label{eq:b3}
\end{equation}
the form of random superposition with learned instead of random phases, and phase-fixed weighted GS runs from it. The network is a permutation-equivariant set transformer: Fourier features of the trap coordinates, an embedding, two attention blocks, a head for the phase and a set-dependent chirp coefficient. Training is unsupervised: the loss is the weighted mean of $\log(\sigma^2+0.01)$ of the trap-intensity relative standard deviation $\sigma$ (in percent) over the start and ten unrolled iterations of phase-fixed weighted GS with an early freeze, the start weighing 0.3 and each iteration 1,
\begin{equation}
\mathcal L=\frac1B\sum_{b=1}^{B}\frac{\sum_{u=0}^{U}\omega_u\log\big(\sigma_{b,u}^2+\epsilon\big)}{\sum_{u=0}^{U}\omega_u},\qquad \epsilon=0.01,\quad \omega_0=0.3,\ \omega_{u\ge1}=1,
\label{eq:b3-loss}
\end{equation}
through a differentiable copy of the rule that was checked against phase-fixed weighted GS, so that the network learns a start from which phase-fixed weighted GS converges sooner. It is not a regression onto converged phases, which are not unique (the global phase and the many fixed points of the weighted GS would make the labels ambiguous). The invariance to the global phase is automatic, since the loss reads intensities only. The trap-set generators are lattices, hexagonal sets, rings, random sets with a minimum spacing and three-dimensional stacks; there is no quasi-periodic family, so that geometry of the benchmark is unseen.

\textbf{Pseudocode.} Algorithm~\ref{alg:b3} lists the inference; training minimizes the loss of Eq.~\eqref{eq:b3-loss} offline.

\begin{algorithm}[t]
\caption{Learned warm start for phase-fixed weighted GS: the network predicts one start phase per trap, and phase-fixed weighted GS runs from the superposition hologram built with them.}
\label{alg:b3}
\begin{algorithmic}[1]
\Require traps (at most 1024, no vortex), released network (or none), freeze index $N$, iteration budget $K$
\If{a model is released} \State $\theta\gets$ network output for the trap list (coordinates in bins, depths in Rayleigh ranges, weights)
  \State $\phi^0\gets\arg\sum_m\sqrt{w_m}\,e^{i(\Delta^m+\theta_m)}$ by Eq.~\eqref{eq:b3}
\Else \State $\phi^0\gets$ none (the answer is phase-fixed weighted GS from its own random start, the provenance \texttt{baseline}) \EndIf
\State \Return phase-fixed weighted GS of Algorithm~\ref{alg:m07} from $\phi^0$ with freeze index $N$ and budget $K$ (point form)
\end{algorithmic}
\end{algorithm}

\textbf{Evaluation and result.} The held-out set was written before training: \wsN{} geometries, generated demanding ones plus the benchmark geometries, with seeds recorded in the training record. The model, with its early freeze, needs a median of \wsModelIt{} iterations to reach one percent against \wsRandomIt{} for a random start with the original freeze, a ratio of \wsRatio{} against the required factor of two, and the final efficiency differs by \wsEffDiff{} percent between the arms against the required one. A control with no network and the early freeze needs a median of \wsControlIt{} iterations on the control subset, and \wsControlNever{} percent of its geometries never reach one percent (\wsModelNever{} percent with the network). The early freeze, not the network, shortens the run, so the hypothesis is refuted. An oracle bound explains why: starting from the trap phases of a converged random-start phase-fixed weighted GS solution, the original freeze still needs many iterations on dense lattices, and only a short freeze shortens the run.

\textbf{Parameters.} The module default for the freeze is \defv{warmstart}{nfix}, the value of phase-fixed weighted GS; the model arm in training and in the evaluation uses the early freeze; the form defaults to \texttt{\defv{warmstart}{form}} because three-dimensional traps need it. With no released model the answer is phase-fixed weighted GS from its own random start with the provenance \texttt{baseline}.

\textbf{Acceptance criterion and gate.} The ladder row records a measured negative. The gate asserts that the ratio of the medians matches the stored one and lies in the interval from 1.5 to 2.0, that the final efficiency differs by more than one percent, that the control supports the explanation, and that no released warm-start model file exists; a second test holds a released model, if one is ever released, to the original criterion. The model is not released, and the committed record of the refusal is the evidence. The remaining open question, a label-based model or a change of the freeze rule, is a new design question, not tuning, and it is not claimed. The module is \texttt{b3\_warm\_start}.

\subsection{Aberration-robust hologram}\label{sec:b4}

After a Zernike correction (Section~\ref{sec:m14}) residual aberrations remain and vary across the field. The fourth proposed method keeps the trap spread low over the whole distribution of residuals instead of at the nominal point. The hypothesis as narrowed after a measurement is that minimizing the expected trap-intensity loss over a distribution of residual Zernike aberrations gives a hologram whose trap spread over an ensemble of residuals is lower than that of the nominal phase-fixed weighted GS: the 95th percentile of the relative standard deviation over a 200-sample ensemble is at least 1.25 times below that of phase-fixed weighted GS, with the mean efficiency at most ten percent lower. Robustness to pointing errors is an option that is measured and reported and not claimed. The original claim was robustness to aberrations and to trap displacement together, with a factor of two; with the real implementation it failed, and the claim was narrowed. The gap is that aberrations vary across the field and are corrected with one global set of low-order Zernike modes~\cite{ebadi2021}; the nearest work on holograms robust to such errors is~\cite{machu}, and no source that was read optimizes the hologram for a distribution of residual aberration.

\textbf{Method.} The expected loss over residual Zernike coefficients $a_k\sim\mathcal N(0,s_a^2)$ for Noll modes 4 to 11, in waves, with piston and tilts excluded, is minimized,
\begin{equation}
\min_\phi\ \mathbb E_{a}\,L\big(\phi+2\pi W(a)\big),\qquad L=\frac{\sum_m(I_m-T_m)^2}{\sum_mT_m^2},\qquad T_m=\ell\,e_{\mathrm{nom}}\frac{w_m}{\sum w},
\label{eq:b4}
\end{equation}
with $W$ the wavefront added to the SLM phase as in the total-phase sum of Nogrette et al.~\cite{nogrette}, $e_{\mathrm{nom}}$ the efficiency of the starting hologram and $\ell$ the target level fraction. The aberration is a phase offset, so the gradient is the unperturbed closed form at the shifted phase. The default design replaces the expectation by a deterministic sigma-point set, plus and minus $s_a\sqrt8$ on one mode at a time,
\begin{equation}
\mathbb E_a\,L\approx\frac1{16}\sum_{k=1}^{8}\sum_{\pm}L\big(\phi\pm2\pi s_a\sqrt8\,Z_k\big),
\label{eq:b4-sigma}
\end{equation}
sixteen points of mean zero and covariance exactly $s_a^2I$ over the eight modes, so that the loss has no sampling noise and a Fletcher--Reeves conjugate gradient converges on it. A second design draws random samples that are resampled at every iteration, with Adam, and adds independent trap displacements; it is selected automatically when pointing errors are requested. The warm start is phase-fixed weighted GS itself.

\textbf{Pseudocode.} Algorithm~\ref{alg:b4} lists the two designs.

\begin{algorithm}[t]
\caption{Aberration-robust hologram: the expected loss over residual Zernike aberrations by the sigma-point design (conjugate gradient) or by random draws (Adam).}
\label{alg:b4}
\begin{algorithmic}[1]
\Require traps, prior $s_a$, level fraction $\ell$, warm iterations $W$, iteration budget $K$, design, samples $q$, rate $\mathrm{lr}$, pointing prior
\State $\phi\gets$ $W$ iterations of phase-fixed weighted GS from the random superposition; $\mathrm{level}\gets\ell\times$ efficiency of $\phi$
\If{design is sigma points}
  \State define $F(\phi)=\frac1{16}\sum_{k=1}^{8}\sum_\pm L\big(\phi\pm2\pi s_a\sqrt8\,Z_k\big)$ and its gradient at the target level $\mathrm{level}$
  \State \Return $\phi$ after $K$ conjugate-gradient iterations on $F$ (Algorithm~\ref{alg:m10})
\Else
  \For{$i=1,\dots,K$}
    \State draw $q$ perturbations (Noll 4 to 11 coefficients $\sim\mathcal N(0,s_a^2)$ and independent trap displacements); average the loss and gradient at $\phi+2\pi W(a)$ over the displaced traps
    \State Adam step with rate $\mathrm{lr}$ (as in Algorithm~\ref{alg:m11})
  \EndFor
  \State \Return $\phi$
\EndIf
\end{algorithmic}
\end{algorithm}

\textbf{Parameters.} The prior $s_a$ is \defv{robust}{aberration} waves per mode, a stated choice that is not sourced; the level fraction is \defv{robust}{levelfraction}, a measured trade-off (a higher level gives a lower factor, because the level trades efficiency for robustness); the warm start has \defv{robust}{warm} iterations; the samples per iteration are \defv{robust}{q}.

\textbf{Acceptance criterion and gate.} Over a 200-sample ensemble the 95th percentile of the spread is at least 1.25 times below that of phase-fixed weighted GS with an efficiency at most ten percent lower (Table~\ref{tab:ladder}). The gate runs the $10\times10$ task with the nominal hologram from the in-test weighted GS and the aberration-robust hologram with 40 iterations, evaluates both on an ensemble of a different seed than the design, and checks the two bounds; a second test only asserts that the pointing case runs and returns finite ratios. The measurement of this report is: \acc{robust}. The aberration claim holds on every seed. The pointing claim is not obtained in the acceptance configuration (a ratio of 1.04 over three seeds, 0.85 to 1.05); the spread under displacement has a first-order slope term and a second-order curvature floor, and with a displacement prior alone and 100 steps the slope falls about fivefold and the ratio reaches 1.48 at 0.10 bins (1.44 to 1.50) at 0.91 of the efficiency, what remains being close to the curvature floor that the spot width fixes. The ensemble is drawn from the same prior as the design, so robustness to a different distribution or to a measured instrument is not established. In the aberrated case C09 the mean spread of the aberration-robust hologram is below that of phase-fixed weighted GS by a smaller factor than the ensemble claim, from a single realization per seed.

\textbf{Complexity.} Each step costs the number of samples (or sigma points) times one forward and adjoint evaluation of the point-trap model, that is, that number of weighted-GS iterations, and the ensemble evaluation is 200 forward evaluations; the method is baked, and the acceptance test takes about a minute on the development workstation.

\textbf{Limits.} The prior is a stated parameter. The benefit depends on the target level and the budget. The model is aberrations and nothing else on the ideal point-trap model. On ideal cases the aberration-robust hologram's spread is a small nominal spread, the price of robustness. The module is \texttt{b4\_robust}.

\subsection{Efficiency--uniformity Pareto front}\label{sec:b5}

Every method trades efficiency against uniformity and reports single points of that trade-off: Di Leonardo et al. report the direct search and the adaptive form at stated weights~\cite{dileonardo}, Polin et al. use one scalar weight~\cite{polin}, and MRAF has one mixing parameter. The fifth proposed method computes the whole front on the same task. The hypothesis is that a continuation over a trade-off parameter $\lambda_{\mathrm{EU}}\in[0,1]$ of one loss, warm-started from the neighbouring value, gives a front whose hypervolume in the efficiency--uniformity plane is at least that of the union of the single points of GS, GAA, direct search and weighted GS, with at least one front point dominating each of GS and GAA, in under a minute on a processor for the $10\times10$ task. A risk anticipated in the design of the method is that the weighted-GS point is already near the corner, so that the front may be flat and its useful region the large or device-affected cases.

\textbf{Method.} The loss over the continuation is
\begin{equation}
L_{\lambda_{\mathrm{EU}}}(\phi)=(1-\lambda_{\mathrm{EU}})\,\frac{\sum_m(I_m-T_m)^2}{\sum_mT_m^2}-\lambda_{\mathrm{EU}}\,\frac{\sum_mI_m}{e_0},\qquad T_m=e_0\frac{w_m}{\sum_kw_k},
\label{eq:b5}
\end{equation}
with $I_m=|V_m|^2$ and $e_0$ the efficiency of the weighted-GS hologram at the uniform end. DinHot divides the two terms by $\sum T^2$ and $e_0$ so that both are dimensionless and of order one, which makes the weight $\lambda_{\mathrm{EU}}$ comparable across tasks; the scaling is part of the method's definition. The start at $\lambda_{\mathrm{EU}}=0$ is phase-fixed weighted GS refined by the same conjugate gradient, and each next $\lambda_{\mathrm{EU}}$ is warm-started from its neighbour and solved with Fletcher--Reeves conjugate gradient. The gradient is the closed form $G_m=[2(1-\lambda_{\mathrm{EU}})(I_m-T_m)/\sum T^2-\lambda_{\mathrm{EU}}/e_0]V_m$ pulled back by the point-trap adjoint. The metric of the front is the hypervolume, the area of the union of the rectangles $[0,p]$ over the points $p=(e,u)$, with the reference point at the origin,
\begin{equation}
\mathrm{HV}=\sum_{i}\big(e_{(i)}-e_{(i-1)}\big)\,\max_{k\ge i}u_{(k)},\qquad e_{(0)}=0,
\label{eq:b5-hv}
\end{equation}
over the non-dominated points sorted by efficiency; the dominance relation keeps the rows for which both columns are maximized. The knee returned as the hologram is the point of largest $e\cdot u$.

\textbf{Pseudocode.} Algorithm~\ref{alg:b5} lists the continuation.

\begin{algorithm}[t]
\caption{Efficiency--uniformity front by continuation over $\lambda_{\mathrm{EU}}$, warm-started from the neighbouring value.}
\label{alg:b5}
\begin{algorithmic}[1]
\Require traps, number of values $n_{\mathrm{EU}}$, warm iterations $W$, conjugate-gradient iterations $K$ per value
\State $\phi\gets$ $W$ iterations of phase-fixed weighted GS from the random superposition; $e_0\gets$ efficiency of $\phi$; $T_m\gets e_0w_m/\sum_kw_k$
\For{$\lambda_{\mathrm{EU}}\in\{0,\frac1{n_{\mathrm{EU}}-1},\dots,1\}$}
  \State $\phi\gets$ $K$ Fletcher--Reeves conjugate-gradient iterations on $L_{\lambda_{\mathrm{EU}}}$ of Eq.~\eqref{eq:b5} from $\phi$ (Algorithm~\ref{alg:m10})
  \State store $\phi$ and the point $(e,u,\sigma)$ of its trap fields
\EndFor
\State non-dominated points $\gets$ the stored points that no other point beats in both $e$ and $u$; $\mathrm{HV}\gets$ Eq.~\eqref{eq:b5-hv}
\State \Return the stored hologram of largest $e\cdot u$ (the knee), the front and $\mathrm{HV}$
\end{algorithmic}
\end{algorithm}

\textbf{Parameters.} The number of values of $\lambda_{\mathrm{EU}}$ is \defv{paretofront}{nlambda}, within the range of ten to twenty set in the design of the method; the warm start is \defv{paretofront}{warm} phase-fixed iterations; the conjugate-gradient iterations per $\lambda_{\mathrm{EU}}$ are the budget (15 in the acceptance test). The ideal device is assumed, so the front does not apply to a device-limited case.

\textbf{Acceptance criterion and gate.} The hypervolume is at least that of the plain GS, the adaptive form, direct search and weighted GS points on the $10\times10$ task (Table~\ref{tab:ladder}). The gate evaluates the baseline points by the metrics, runs the Pareto front with 20 values of $\lambda_{\mathrm{EU}}$, and requires the hypervolume bound, at least one front point dominating each of plain GS and the adaptive form, and a run time under a minute. The measurement of this report is: \acc{paretofront}. The efficiency changes by only a few percent across the whole front while the uniformity falls from one to a small value, so on this task the trade-off is steep and the useful part of the front is the high-uniformity end, which is the flat-front risk anticipated in the design. In the benchmark the stored front, the converged run of the first seed, reaches at least the hypervolume of the published points on \benchFivePass{} of the \benchFiveN{} geometries with a front against their medians at 100 iterations on the ideal device, and on \benchFivePassThirty{} against their medians at 30; at 100 iterations it falls short on \benchFiveLost, by at most \benchFiveMargin{}. The count at 100 iterations moves between runs, because the medians of a cell depend on the seeds its wall-clock bound admits.

\textbf{Complexity.} The method is $n_{\mathrm{EU}}$ warm-started conjugate-gradient runs of the budget's iterations each, with one forward and one adjoint evaluation per line-search step; the front of the acceptance task takes a fraction of a minute on the development workstation, under the limit of one minute, and the time depends on the machine.

\textbf{Limits.} The criterion is one-sided and the margin is small: the test cannot distinguish a front that is the same as a single good point. The hypervolume depends on the reference point and is dominated by the high-uniformity region. The front is for one task and one setup. The module is \texttt{b5\_pareto}, whose hypervolume and dominance functions are public for the benchmark.
